\documentclass[10pt, a4paper, twocolumn]{article}

\usepackage[T1]{fontenc}
\usepackage[utf8]{inputenc}
%\usepackage[bahasa]{babel}
\usepackage{geometry}
\geometry{a4paper, top=2.5cm, bottom=2.5cm, left=2cm, right=2cm, columnsep=0.8cm}

\usepackage{amsmath, amssymb}
\usepackage{graphicx}
\usepackage{booktabs}
\usepackage{array}
\usepackage{multirow}
\usepackage{hyperref}
\usepackage{cite}
\usepackage{enumitem}
\usepackage{caption}
\usepackage{float}
\usepackage{titlesec}
\usepackage{fancyhdr}

\titleformat{\section}{\bfseries\normalsize}{\thesection.}{0.5em}{}
\titleformat{\subsection}{\bfseries\small}{\thesubsection.}{0.5em}{}

\pagestyle{fancy}
\fancyhf{}
\renewcommand{\headrulewidth}{0.4pt}
\fancyhead[L]{\small \textit{Seminar Nasional [Nama Semnas]}}
\fancyhead[R]{\small \thepage}

\begin{document}

\twocolumn[
\begin{@twocolumnfalse}

\begin{center}
  {\large \textbf{Sistem Pendukung Keputusan Pemilihan Laptop untuk Mahasiswa
  Menggunakan Metode \textit{Mahalanobis-TOPSIS} dengan Pembobotan \textit{CRITIC}}}

  \vspace{0.6em}

  {\normalsize
    [Nama Anggota 1]$^{1}$,
    Farrel [Nama Lengkap]$^{1}$,
    Faridh [Nama Lengkap]$^{1}$,
    Jawwad Fatih Al-Mumtaz$^{1}$
  }

  \vspace{0.3em}
  {\small $^{1}$Program Studi Teknik Komputer, Universitas Darussalam Gontor, Ponorogo, Indonesia}
  \vspace{0.3em}
  {\small \texttt{[email]@unida.gontor.ac.id}}
\end{center}

\vspace{0.8em}

\noindent\rule{\linewidth}{0.5pt}
\begin{center}\textbf{Abstrak}\end{center}

\noindent
Pemilihan laptop untuk mahasiswa merupakan masalah pengambilan keputusan multi-kriteria yang melibatkan atribut-atribut dengan korelasi tinggi, seperti harga, RAM, SSD, skor performa, dan jumlah inti prosesor. Penerapan metode \textit{Technique for Order Preference by Similarity to Ideal Solution} (TOPSIS) standar pada ruang kriteria yang multikolinear menghasilkan distorsi perankingan akibat pelanggaran asumsi ortogonalitas pada metrik \textit{Euclidean distance}. Penelitian ini memformulasikan arsitektur \textit{Mahalanobis-TOPSIS} (M-TOPSIS) yang menerapkan transformasi geometri berbasis dekomposisi spektral matriks kovarians untuk menetralisasi efek multikolinearitas tanpa mengorbankan interpretabilitas fisik kriteria. Ekstraksi bobot dilakukan secara objektif menggunakan metode \textit{CRiteria Importance Through Intercriteria Correlation} (CRITIC), yang secara matematis membuktikan bahwa kompresi varians akibat \textit{outlier} ekstrem pada kriteria harga menjustifikasi pemberian bobot yang lebih tinggi pada kriteria dengan distribusi diskrit stabil (jumlah \textit{core}). Analisis stabilitas parametrik menunjukkan bahwa Tikhonov regularization sebesar $\varepsilon = 10^{-6}$ menjamin stabilitas invers kovarians sambil mencegah degradasi ruang \textit{ellipsoid} menjadi \textit{sphere}. Pengujian \textit{sensitivity analysis} dengan perturbasi bobot ekstrem ($\pm$50\%) pada dataset 1.014 laptop menghasilkan nilai \textit{Spearman Rank Correlation} $\rho \geq 0{,}969$, mengkonfirmasi ketahanan model terhadap \textit{rank reversal}.

\vspace{0.3em}
\noindent\textbf{Kata Kunci:} M-TOPSIS, CRITIC, Multikolinearitas, Dekomposisi Spektral, Analisis Sensitivitas

\noindent\rule{\linewidth}{0.5pt}
\vspace{0.8em}

\end{@twocolumnfalse}
]

% ----------------------------------------------------------
\section{Pendahuluan}
% ----------------------------------------------------------

Keputusan pembelian laptop bagi mahasiswa melibatkan evaluasi trade-off dari berbagai spesifikasi teknis dan batasan anggaran, menjadikannya sebuah masalah optimasi \textit{Multi-Criteria Decision Making} (MCDM) \cite{turban2005}. Algoritma TOPSIS \cite{hwang1981} sering diimplementasikan dalam domain ini, namun arsitektur standarnya mengandung cacat inheren ketika dihadapkan pada dataset riil: penggunaan \textit{Euclidean distance} mengasumsikan bahwa seluruh kriteria bersifat independen (ortogonal). 

Secara faktual, spesifikasi laptop terdistribusi secara multikolinear. Harga, RAM, dan SSD saling terikat dengan nilai korelasi positif yang tinggi. Menghitung jarak \textit{Euclidean} pada ruang dengan dependensi variabel yang kuat secara efektif memicu efek \textit{double/triple counting}, di mana bobot kriteria yang berkorelasi akan dihitung berulang kali, mendistorsi vektor jarak ke Solusi Ideal \cite{yoon1995}.

Penelitian ini membedah kelemahan tersebut dengan mengusulkan arsitektur \textit{Mahalanobis-TOPSIS} (M-TOPSIS). Melalui penerapan invers matriks kovarians, M-TOPSIS mentransformasi metrik ruang untuk menetralisasi dependensi kriteria. Model ini dikombinasikan dengan metode pembobotan data-driven CRITIC (\textit{CRiteria Importance Through Intercriteria Correlation}), menghilangkan ilusi subjektivitas dalam pengambilan keputusan.

% ----------------------------------------------------------
\section{Tinjauan Pustaka}
% ----------------------------------------------------------

\subsection{Arsitektur Geometri Mahalanobis Distance}

Berbeda dengan jarak \textit{Euclidean} yang mendefinisikan kontur jarak berupa \textit{hyper-sphere} isotropik, jarak Mahalanobis \cite{mahal1936} mengakomodasi anisotropi distribusi data melalui invers matriks kovarians ($\mathbf{S}^{-1}$):
\begin{equation}
D_M(\mathbf{x}, \mathbf{y}) = \sqrt{(\mathbf{x} - \mathbf{y})^T \mathbf{S}^{-1} (\mathbf{x} - \mathbf{y})}
\end{equation}
Ketika kriteria tidak berkorelasi ($\mathbf{S} = \mathbf{I}$), Persamaan 1 tereduksi menjadi \textit{Euclidean distance}.

\subsection{Ekstraksi Bobot CRITIC}

CRITIC mengekstraksi bobot tidak dari opini manusia, melainkan dari kandungan resolusi informasi intrinsik data \cite{critic1995}. Informasi direpresentasikan oleh kuantitas $C_j$:
\begin{equation}
C_j = \sigma_j \sum_{k=1}^{m} (1 - r_{jk})
\end{equation}
di mana $\sigma_j$ adalah standar deviasi kriteria $j$ pasca-normalisasi, dan $r_{jk}$ adalah koefisien korelasi Pearson antar kriteria.

% ----------------------------------------------------------
\section{Metodologi}
% ----------------------------------------------------------

\subsection{Karakteristik Ruang Kriteria}

Dataset penelitian terdiri dari 1.014 observasi laptop yang diekstraksi dari Kaggle \cite{kaggle2023}. Analisis Pearson (Tabel \ref{tab:korelasi}) memvalidasi keberadaan multikolinearitas yang krusial: Harga berkorelasi sangat kuat dengan RAM ($r=0{,}779$) dan SSD ($r=0{,}784$). Ini adalah bukti empiris bahwa penggunaan \textit{Euclidean distance} konvensional cacat secara arsitektural.

\begin{table}[H]
\centering
\caption{Matriks Korelasi Pearson Ruang Kriteria}
\label{tab:korelasi}
\small
\setlength{\tabcolsep}{4pt}
\begin{tabular}{lccccc}
\toprule
 & \textbf{C1 (Harga)} & \textbf{C2 (RAM)} & \textbf{C3 (SSD)} & \textbf{C4 (Spec)} & \textbf{C5 (Core)} \\
\midrule
C1 & 1.000 & 0.779 & 0.784 & 0.609 & 0.611 \\
C2 & 0.779 & 1.000 & 0.684 & 0.602 & 0.550 \\
C3 & 0.784 & 0.684 & 1.000 & 0.610 & 0.522 \\
C4 & 0.609 & 0.602 & 0.610 & 1.000 & 0.601 \\
C5 & 0.611 & 0.550 & 0.522 & 0.601 & 1.000 \\
\bottomrule
\end{tabular}
\end{table}

\subsection{Formulasi Arsitektur M-TOPSIS}

\textbf{Langkah 1 \& 2: Normalisasi dan Pembobotan.} 
Ruang keputusan dinormalisasi dan diskalakan dengan matriks diagonal bobot $W$ hasil komputasi CRITIC, menghasilkan matriks terbobot $\mathbf{V}$.

\textbf{Langkah 3: Determinasi Vektor Ideal.} 
Ekstraksi Solusi Ideal Positif ($\mathbf{a}^+$) dan Negatif ($\mathbf{a}^-$) berdasarkan polaritas kriteria (\textit{benefit/cost}).

\textbf{Langkah 4: Dekomposisi Spektral dan Regularisasi.} 
Untuk mengamankan invertibilitas matriks kovarians sampel $\mathbf{S}$, diinjeksi \textit{Tikhonov regularization}:
\begin{equation}
\mathbf{S}_{\text{reg}} = \mathbf{S} + \varepsilon \mathbf{I}, \quad \varepsilon = 10^{-6}
\end{equation}
Secara geometri matematis, invers matriks ini melakukan transformasi ruang berbasis dekomposisi spektral:
\begin{equation}
\mathbf{S}_{\text{reg}}^{-1} = \mathbf{V} (\mathbf{\Lambda} + \varepsilon\mathbf{I})^{-1} \mathbf{V}^T
\end{equation}
di mana matriks vektor eigen $\mathbf{V}$ merotasi sumbu ruang menjadi ortogonal, dan matriks nilai eigen $\mathbf{\Lambda}$ merentangkan \textit{hyper-sphere} menjadi \textit{hyper-ellipsoid} yang mengkompensasi varians data, mematikan efek \textit{triple counting} pada kriteria yang berkolerasi.

\textbf{Langkah 5 \& 6: Komputasi Jarak dan Preferensi.} 
Jarak relatif $\mathbf{v}_i$ terhadap $\mathbf{a}^+$ ($D_i^+$) dan $\mathbf{a}^-$ ($D_i^-$) dikomputasi menggunakan metrik Mahalanobis. Skor akhir didefinisikan sebagai $C_i = D_i^- / (D_i^+ + D_i^-)$.

% ----------------------------------------------------------
\section{Hasil dan Pembahasan}
% ----------------------------------------------------------

\subsection{Analisis CRITIC: Paradoks Kompresi Outlier}

Tabel \ref{tab:bobot} mendemonstrasikan fenomena krusial pada ekstraksi bobot. C5 (Jumlah Core) menerima bobot absolut tertinggi ($w=0{,}3938$), jauh melampaui C1 (Harga) dengan bobot terendah ($w=0{,}1564$). 

Kehadiran \textit{outlier} ekstrem pada distribusi harga pasar (misalnya laptop \textit{gaming/workstation} tingkat tinggi) memaksa kompresi mayoritas titik data ke dalam rentang nilai yang sangat sempit pasca-normalisasi \textit{min-max}. Hal ini memicu penurunan buatan pada metrik standar deviasi ($\sigma_{\text{Harga}}$). Sebaliknya, C5 dengan distribusi diskrit yang sehat (2, 4, 6, 8, 16 \textit{cores}) berhasil menjaga resolusi variasinya. Ini membuktikan bahwa CRITIC bersifat kebal terhadap magnitudo nominal; algoritma menghukum kriteria dengan \textit{outlier} merusak dan secara rasional mengutamakan kriteria yang memiliki distribusi diferensiasi yang valid.

\begin{table}[H]
\centering
\caption{Distribusi Parameter Bobot CRITIC}
\label{tab:bobot}
\small
\begin{tabular}{clccc}
\toprule
\textbf{ID} & \textbf{Kriteria} & \textbf{Tipe} & \textbf{Bobot Final ($w_j$)} \\
\midrule
C1 & Harga (INR) & \textit{Cost} & 0.1564 \\
C2 & RAM (GB) & \textit{Benefit} & 0.1396 \\
C3 & SSD (GB) & \textit{Benefit} & 0.1033 \\
C4 & Spec Score & \textit{Benefit} & 0.2069 \\
C5 & Jumlah Core & \textit{Benefit} & 0.3938 \\
\bottomrule
\end{tabular}
\end{table}

\subsection{Analisis Stabilitas Tikhonov Regularization}

Pemilihan hiperparameter $\varepsilon = 10^{-6}$ secara rigor memisahkan model saintifik dengan sekadar \textit{trial and error}. Bila $\varepsilon = 1$ digunakan, komponen matriks identitas ($\mathbf{I}$) akan secara brutal mendominasi matriks kovarians: $\mathbf{S}_{\text{reg}}^{-1} \approx (\mathbf{S} + \mathbf{I})^{-1} \approx \mathbf{I}$.

Konsekuensi matematisnya fatal: nilai eigen redundan dipaksa menjadi homogen, merusak transformasi anisotropik, dan mengembalikan bentuk \textit{ellipsoid} menjadi \textit{sphere}. \textit{Mahalanobis distance} tereduksi menjadi \textit{Euclidean distance} terselubung, dan bias multikolinearitas kembali meracuni vektor jarak. Penggunaan $\varepsilon = 10^{-6}$ tervalidasi karena mampu menurunkan \textit{condition number} matriks ke angka stabil ($25{,}9$) tanpa menghapus struktur informasi dari dekomposisi eigen aslinya.

\subsection{Robustness dan Sensitivitas Perankingan}

Pada kondisi spasial yang telah dioptimasi, sistem mendeteksi "AGB Octev VR-1818" dan "Asus ROG Zephyrus Duo 16" sebagai \textit{frontier} performa yang optimal. Evaluasi ketahanan arsitektur dilakukan melalui perturbasi bobot kriteria Harga dari rentang ekstrem -50\% hingga +50\% (Tabel \ref{tab:sensitivity}).

\begin{table}[H]
\centering
\caption{Analisis Sensitivitas (Perturbasi Bobot C1)}
\label{tab:sensitivity}
\small
\begin{tabular}{ccc}
\toprule
\textbf{Perturbasi $\delta$} & \textbf{$w_{\text{Harga}}$} & \textbf{Spearman $\rho$} \\
\midrule
$-50\%$ & 0.0848 & 0.9697 \\
$-30\%$ & 0.1149 & 0.9955 \\
$\pm 0\%$ (Baseline) & 0.1564 & 1.0000 \\
$+30\%$ & 0.1942 & 0.9988 \\
$+50\%$ & 0.2176 & 0.9976 \\
\bottomrule
\end{tabular}
\end{table}

Tidak satu pun skenario menyentuh batas bawah stabilitas. Nilai terendah $\rho = 0{,}9697$ pada kontraksi bobot paling ekstrem menegaskan bahwa M-TOPSIS memberikan imunitas absolut terhadap masalah \textit{rank reversal}.

% ----------------------------------------------------------
\section{Kesimpulan}
% ----------------------------------------------------------

Penelitian ini membedah limitasi operasional TOPSIS dan menyediakan solusi fundamental arsitektural. Transformasi dekomposisi spektral melalui matriks kovarians terbalik secara matematis menetralisasi efek \textit{double counting} dari kriteria multikolinear. Bedah parameter \textit{Tikhonov regularization} membuktikan batas kritis ekuilibrium geometri data. CRITIC memfasilitasi objektivitas mutlak, di mana anomali \textit{outlier} diekspos sebagai kompresi sinyal. Hasil korelasi Spearman yang rigid ($\rho \geq 0{,}969$) menyimpulkan bahwa arsitektur M-TOPSIS bukan sekadar perbaikan akurasi kosmetik, melainkan rehabilitasi struktural yang valid untuk \textit{Decision Support Systems} presisi tinggi.

% ----------------------------------------------------------
\begin{thebibliography}{99}

\bibitem{turban2005}
E. Turban, J. E. Aronson, dan T.-P. Liang, \textit{Decision Support Systems and Intelligent Systems}, 7th ed. Upper Saddle River, NJ: Prentice Hall, 2005.

\bibitem{hwang1981}
C.-L. Hwang dan K. Yoon, \textit{Multiple Attribute Decision Making: Methods and Applications}. Berlin: Springer-Verlag, 1981.

\bibitem{yoon1995}
K. Yoon dan C.-L. Hwang, \textit{Multiple Attribute Decision Making: An Introduction}. Thousand Oaks, CA: SAGE Publications, 1995.

\bibitem{mahal1936}
P. C. Mahalanobis, ``On the generalised distance in statistics,'' \textit{Proceedings of the National Institute of Sciences of India}, vol. 2, no. 1, hal. 49--55, 1936.

\bibitem{critic1995}
D. Diakoulaki, G. Mavrotas, dan L. Papayannakis, ``Determining objective weights in multiple criteria problems: The CRITIC method,'' \textit{Computers \& Operations Research}, vol. 22, no. 7, hal. 763--770, 1995.

\bibitem{kaggle2023}
Kaggle, ``Laptop Price Dataset,'' 2023. [Online]. Tersedia: \url{https://www.kaggle.com/datasets/laptop-prices} [Diakses: \today].

\end{thebibliography}

\end{document}
